% !PRINTED 402
\weberpage{402}
%<<CONT

4. \gesperrt{Ist $\mathfrak{a}$ ein zu $Q$ teilerfremdes Ideal, so liefert uns
jedes den Bedingungen (12) genügende Zahlenpaar $x_1, x_2$
eine durch $\mathfrak{a}$ teilbare positive Zahl $y$ mit positiver Norm
der Ordnung $[Q]$}

\weq{13}{y = a\,x_1 + \frac{b+\sqrt{D}}{2}x_2,}

\gesperrt{und unter diesen Zahlen sind keine zwei assoziiert.}

\begin{center}
\S\,111. \gesperrt{Die Dirichletsche Grenzformel.}
\end{center}

Bedeutet $t$ eine positive Konstante, so ist durch die Bedingung:

\weq{1}{N(y) = a\,(a\,x_1^2 + b\,x_1 x_2 + c\,x_2^2) \lessgtr t,} % Source prints a combined less/greater relation; \lessgtr is the matching corpus macro.
und, bei positiver Diskriminante, durch die Isolierungsbedingungen,
in einer Ebene, in der $x_1, x_2$ rechtwinkelige Koordinaten sind,
ein Gebiet $F_t$ begrenzt, das bei negativer Diskriminante durch
eine Ellipse, bei positiver Diskriminante durch einen Hyperbelbogen
 und durch gerade Linien begrenzt ist. Die Punkte, deren
Koordinaten $x_1, x_2$ ganze rationale Zahlen sind, heißen \gesperrt{Gitterpunkte}.

Einem ganzen System von assoziierten Zahlen, deren Norm
positiv und kleiner als $t$ ist, entspricht dann bei positiver Diskriminante
 ein Gitterpunkt in $F_t$ und bei negativer Diskriminante
zwei und in den beiden Ausnahmefällen vier und sechs Gitterpunkte.

Bezeichnet $Z_t$ die Anzahl der Gitterpunkte in dem Gebiete
$F_t$ und $V$ die Fläche des Gebietes $F_t$, so ist [Bd. II, § 194, (6)]:

\weq{2}{V = Z_t\,t^{-1} + R_t\,t^{-\frac{1}{2}},}
worin $R_t$ mit unendlich wachsendem $t$ nicht unendlich wird.

Bei negativer Diskriminante erhält man $V$ aus dem bekannten
Flächeninhalt der Ellipse:

\weq{3}{V = \frac{2\pi}{a\sqrt{-D}},}
und bei positiver Diskriminante erhält man ihn am einfachsten,
indem man nach § 110, (7) an Stelle von $x_1, x_2$ die Integrationsvariablen
 $\xi_1, \xi_2$ einführt, deren Grenzen $0, 1$ und $-\infty, 0$ sind:

\weq{4}{V = \frac{\log\varepsilon}{a\sqrt{D}}.}
